\documentclass[twocolumn,prl,aps]{revtex4-2}
\usepackage{graphicx}
\usepackage{amsmath}

\begin{document}

\title{Automated Discovery of a {{BH_TYPE}} Black Hole in a 4D Spacetime Simulation}

\author{QuantumVerse AI Agent}
\affiliation{QuantumVerse Simulator, Autonomous Discovery Engine}

\begin{abstract}
{{LLM_ABSTRACT}}
\end{abstract}

\maketitle

\section{Introduction}
The QuantumVerse Simulator utilizes a PPO reinforcement learning agent to explore the parameter space of General Relativity. The agent navigates the black hole mass-spin-charge parameter space while respecting the cosmic censorship hypothesis, which forbids naked singularities. Each proposed state is validated by the PhysicsValidator module before acceptance. The discovered object is a {{BH_TYPE}} black hole.

\section{Methodology}
The RL agent explored the parameter space near the {{BH_TYPE}} limit ($a^2 + Q^2 \approx M^2$ for Kerr; Schwarzschild by construction when spin = 0). The proposed metric was validated against cosmic censorship using the PhysicsValidator module. The final validated parameters were determined to be:
\begin{itemize}
    \item Mass ($M$): {{MASS}} $M_\odot$
    \item Spin ($a$): {{SPIN}}
    \item Charge ($Q$): {{CHARGE}}
\end{itemize}

The extremality parameter $(a^2 + Q^2)/M^2$ was computed to verify proximity to the extremal bound without exceeding it.

\section{Results}
The spacetime curvature was evaluated at the probe location using the {{BH_TYPE}} metric curvature invariants. The Kretschmann scalar was found to be $K = {{KRETSCHMANN}}$. The gravitational redshift at the probe radius is $z = {{REDSHIFT}}$. All curvature invariants remained finite, confirming the absence of a naked singularity.

\section{Conclusion}
The AI agent successfully navigated the parameter space to find a stable {{BH_TYPE}} black hole without violating cosmic censorship. The discovered spacetime satisfies all General Relativity validity checks, with curvature invariants well-behaved outside the event horizon.

\end{document}
